%==================================================================
% HALAMAN JUDUL (sampul dalam)
% Acuan: Pedoman Lampiran 5
%==================================================================
\begin{titlepage}
\thispagestyle{empty}%
\centering

\vspace*{0.3cm}

%--- Judul: 14 pt, tebal, kapital, spasi 1,5 -----------------------
\begin{spacing}{1}
\begin{center}
{\fontsize{14}{16.8}\selectfont\bfseries \MakeUppercase{\judulid}\par}
\end{center}
\end{spacing}

\vspace{2\baselineskip}

%--- Jenis dokumen: 12 pt, tebal, spasi 1,5 -----------------------
\begin{spacing}{1}
\begin{center}
{\fontsize{12}{14.4}\selectfont\bfseries \juduljenis\par}
\end{center}
\end{spacing}

\vspace{2\baselineskip}

\includegraphics[width=4cm,height=4cm,keepaspectratio]{logo-poliwangi.png}

\vspace{2\baselineskip}

%--- Kalimat pengantar, spasi 1,5 ---------------------------------
\begin{spacing}{1}
\begin{center}
{\fontsize{12}{14.4}\selectfont
 \ifthenelse{\equal{\Jenis}{Proposal}}{%
   Proposal Tugas Akhir ini dibuat dan diajukan sebagai salah satu
   syarat untuk melaksanakan Tugas Akhir pada Program Studi
   \jenjang\ \prodi\ \institusi.\par
 }{%
   Tugas Akhir ini dibuat dan diajukan untuk memenuhi salah satu
   syarat kelulusan Program Studi \jenjang\ \prodi\ dan mencapai
   gelar \gelar.\par
 }}
\end{center}
\end{spacing}

\vspace{2\baselineskip}

%--- Nama dan NIM: 12 pt, tebal, spasi 1,5 ------------------------
\begin{spacing}{1}
\begin{center}
{\fontsize{12}{14.4}\selectfont\bfseries
 Oleh:\par
 \penulis\par
 NIM. \nim\par}
\end{center}
\end{spacing}

\vspace{2\baselineskip}

%--- Program studi dan tahun: 16 pt, tebal ------------------------
\begin{spacing}{1}
\begin{center}
{\fontsize{16}{19.2}\selectfont\bfseries
 PROGRAM STUDI \MakeUppercase{\prodi}\par
 \MakeUppercase{\institusi}\par
 \tahun\par}
\end{center}
\end{spacing}

\end{titlepage}
